\section{Introdução}

O objetivo da prática foi comparar processamento em CPU e GPU em dois
problemas: a transformada rápida de Fourier (FFT) e o conjunto de Mandelbrot.
Foi autorizado o uso de Python em substituição ao MATLAB. Assim, NumPy foi
usado na CPU, CuPy representou operações vetorizadas sobre dados residentes
na GPU e Numba CUDA foi usado para implementar um kernel explícito.

Os experimentos foram executados em uma NVIDIA GeForce RTX 2060 com 6 GB de
memória. Cada tempo apresentado corresponde à mediana de cinco execuções.
